%%
%% This is file `sample-sigconf-authordraft.tex',
%% generated with the docstrip utility.
%%
%% The original source files were:
%%
%% samples.dtx  (with options: `all,proceedings,bibtex,authordraft')
%% 
%% IMPORTANT NOTICE:
%% 
%% For the copyright see the source file.
%% 
%% Any modified versions of this file must be renamed
%% with new filenames distinct from sample-sigconf-authordraft.tex.
%% 
%% For distribution of the original source see the terms
%% for copying and modification in the file samples.dtx.
%% 
%% This generated file may be distributed as long as the
%% original source files, as listed above, are part of the
%% same distribution. (The sources need not necessarily be
%% in the same archive or directory.)
%%
%%
%% Commands for TeXCount
%TC:macro \cite [option:text,text]
%TC:macro \citep [option:text,text]
%TC:macro \citet [option:text,text]
%TC:envir table 0 1
%TC:envir table* 0 1
%TC:envir tabular [ignore] word
%TC:envir displaymath 0 word
%TC:envir math 0 word
%TC:envir comment 0 0
%%
%% The first command in your LaTeX source must be the \documentclass
%% command.
%%
%% For submission and review of your manuscript please change the
%% command to \documentclass[manuscript, screen, review]{acmart}.
%%
%% When submitting camera ready or to TAPS, please change the command
%% to \documentclass[sigconf]{acmart} or whichever template is required
%% for your publication.
%%
%%
\documentclass[sigplan,10pt,review,anonymous]{acmart}
\renewcommand\footnotetextcopyrightpermission[1]{}
\settopmatter{
  printfolios=true,
  printacmref=false,
  printccs=false
}

\usepackage{graphicx}
\usepackage{booktabs}
\usepackage{algorithm}
\usepackage{algpseudocode}
\usepackage{amsmath}
\usepackage{subcaption}
\usepackage{tikz}
\usetikzlibrary{arrows.meta,positioning,calc}


% 把我给你的 4 行代码粘贴在这里！！！
\exhyphenpenalty=0
\hyphenpenalty=0
\tolerance=500
\emergencystretch=3em

%%
%% \BibTeX command to typeset BibTeX logo in the docs
\AtBeginDocument{%
  \providecommand\BibTeX{{%
    Bib\TeX}}}

\setcopyright{none}


\begin{document}
% \pagestyle{plain}

%%
%% The "title" command has an optional parameter,
%% allowing the author to define a "short title" to be used in page headers.
\title{KVCore: Sparsity-Guided KV Cache Lifecycle Management for Long-Context LLM Serving}

%%
%% The "author" command and its associated commands are used to define
%% the authors and their affiliations.
%% Of note is the shared affiliation of the first two authors, and the
%% "authornote" and "authornotemark" commands
%% used to denote shared contribution to the research.
\author{Anonymous Author(s)}
\renewcommand{\shortauthors}{Anonymous et al.}

%%
%%
%% The abstract is a short summary of the work to be presented in the
%% article.
\begin{abstract}

Modern LLM serving frameworks commonly manage KV cache at block granularity to support dynamic allocation, reduce memory fragmentation, and enable continuous batching. However, block-based management does not fundamentally relieve KV storage pressure during decoding. Since request output lengths are unknown, schedulers that keep GPU KV memory highly utilized may admit batches that initially fit the current budget but continue generating longer than expected. As active requests keep appending new KV blocks, the remaining GPU memory can become insufficient for the next decode iteration, forcing the runtime to preempt requests and reclaim KV space, which increases latency and destabilizes serving.

We propose KVCore, a sparsity-guided KV lifecycle management system. KVCore builds on the fact that KV cache is inherently sparse: each decode step can often preserve generation quality by accessing only a small subset of historical KV entries. If the runtime knew the future minimal KV subsets, it could optimally decide which KV blocks to keep, offload, prefetch, or evict. Since this oracle is unavailable, KVCore uses practical attention-derived sparse scores to approximate near-future minimal KV sets and turns them into block-level lifecycle hints. By aligning sparse scoring with block-based KV management, KVCore proactively reduces GPU KV pressure, mitigates request preemption, and improves end-to-end serving throughput by 1.5--2.3$\times$ while preserving generation quality.
\end{abstract}

%%
%% The code below is generated by the tool at http://dl.acm.org/ccs.cfm.
%% Please copy and paste the code instead of the example below.
%%
% \begin{CCSXML}
% <ccs2012>
%    <concept>
%        <concept_id>10010520.10010521</concept_id>
%        <concept_desc>Computer systems organization~Architectures</concept_desc>
%        <concept_significance>500</concept_significance>
%        </concept>
%    <concept>
%        <concept_id>10010147.10010178</concept_id>
%        <concept_desc>Computing methodologies~Artificial intelligence</concept_desc>
%        <concept_significance>500</concept_significance>
%        </concept>
%    <concept>
%        <concept_id>10002951.10003152.10003520</concept_id>
%        <concept_desc>Information systems~Storage management</concept_desc>
%        <concept_significance>500</concept_significance>
%        </concept>
%  </ccs2012>
% \end{CCSXML}
% CCS concepts and keywords are omitted in the anonymous submission.

% \ccsdesc[500]{Computer systems organization~Architectures}
% \ccsdesc[500]{Computing methodologies~Artificial intelligence}
% \ccsdesc[500]{Information systems~Storage management}

%%
%% Keywords. The author(s) should pick words that accurately describe
%% the work being presented. Separate the keywords with commas.
\keywords{KV Cache Management, Sparse Attention, KV Cache Offloading, Paged Attention}

%%
%% This command processes the author and affiliation and title
%% information and builds the first part of the formatted document.
\maketitle

\section{Introduction}
\label{sec:intro}

KV cache is one of the core runtime states in modern LLM serving systems.
During autoregressive inference, each generated token appends new key and value states to every transformer layer, and future tokens reuse these historical states for attention computation.
Unlike model weights, whose footprint is fixed after loading, KV cache is request-dependent and grows with both the number of active requests and their generated lengths.
Modern LLM servers therefore manage KV cache at block granularity: as shown in \autoref{fig:kvcache}, each request's logical KV sequence is divided into fixed-size blocks and mapped to physical blocks from a global GPU KV block pool.
This block-based abstraction avoids large contiguous allocations, reduces fragmentation, and enables dynamic allocation, reclamation, and continuous batching.
However, it only changes how KV memory is organized; active requests still keep appending new KV blocks throughout decoding.

This continuously growing KV footprint makes high-utilization scheduling fragile under fixed GPU memory capacity.
To improve batching efficiency and GPU throughput, schedulers typically keep the KV block pool highly utilized, admitting new requests or continuing a running batch as long as the current KV footprint appears feasible for the next decode iteration.
However, request output lengths are unknown at scheduling time.
A batch that fits the current budget may continue decoding for many more steps, gradually consuming the remaining GPU blocks.
When the free KV space is no longer sufficient to support the next decode iteration, the runtime must preempt some active requests, whose KV cache may be swapped, released and later recomputed, or delayed until enough blocks become available.
As shown in \autoref{fig:block-preempt}, this behavior leads to repeated KV pressure cycles and increasing preemption rates under larger batches and longer outputs, raising latency and destabilizing continuous batching.

KV sparsity offers a natural way to reduce this pressure.
Existing KV sparsity methods exploit the fact that not all historical KV entries contribute equally to attention; by retaining or attending to only important entries, they reduce both attention computation and KV memory footprint.
However, most such methods are designed as one-shot compression or per-step sparse selection.
One-shot compression ignores how KV importance evolves during subsequent decoding, while per-step selection only decides which entries participate in the current attention computation.
Moreover, many methods operate at token granularity or rely on attention-kernel-specific information, making their sparse decisions difficult to expose as stable runtime-level signals in block-based serving systems.
What is missing is a low-overhead block-level signal that can track KV importance over time and guide runtime lifecycle decisions.

The key opportunity comes from the inherent sparsity of attention.
For each decode step $i$, there exists an ideal minimal KV set $S_i^\star$ that is sufficient to preserve generation quality.
The future sequence ${S_i^\star}$ forms an ideal sparse-access trace: if known, it would directly tell the runtime which blocks should stay resident, which can be temporarily offloaded, which should be prefetched, and which can be permanently evicted.
Unfortunately, this oracle trace is unavailable because future queries and their exact sparse KV requirements are unknown.
Nevertheless, practical sparse scoring can provide a useful approximation.
As illustrated in \autoref{fig:trace}, KV importance changes across decode steps, but important KV sets are locally stable across nearby steps; sparse eviction also exhibits spatial continuity, allowing token-level sparse evidence to be lifted to block granularity.
These observations suggest that sparse scores can be used not merely as local pruning decisions, but as online hints of near-future KV value.

Based on these observations, we propose KVCore, a sparsity-guided KV lifecycle management system for block-based LLM serving.
KVCore introduces a block-level sparse scoring method that estimates KV block importance using lightweight GPU-resident summaries, avoiding full attention-score materialization and aligning sparse decisions with the runtime's block abstraction.
It then uses these scores as online approximations of near-future sparse access to guide KV lifecycle actions, including GPU residency, CPU offload, prefetching, and permanent eviction.
KVCore implements these actions through block-state management and asynchronous block movement, allowing the runtime to proactively reduce GPU KV pressure before preemption becomes necessary.
Our evaluation shows that KVCore preserves generation quality while improving end-to-end serving throughput by 1.5--2.3$\times$.

This paper makes the following contributions:
\begin{enumerate}
\item \textbf{Problem and abstraction.}
We identify that dynamic KV growth under unknown output lengths destabilizes high-utilization scheduling in block-based LLM serving.
We further formalize the future sequence of minimal sparse KV sets ${S_i^\star}$ as an ideal sparse-access trace for KV lifecycle management.

\item \textbf{Block-level sparse scoring.}
We design a lightweight query-aware scoring method that summarizes KV blocks on GPU and estimates block importance without materializing full attention scores.
The scoring signal is continuously refreshed and directly matches the block abstraction used by modern KV managers.

\item \textbf{Sparsity-guided lifecycle management.}
We build KVCore to translate block-level sparse scores into runtime lifecycle actions, including GPU residency, CPU offload, prefetching, and permanent eviction.
This reduces GPU KV pressure, mitigates request preemption, and improves serving throughput while preserving generation quality.
\end{enumerate}


\begin{figure}[htbp]  % htbp 是浮动位置参数，让图片尽量靠近文字
  \centering   
  \includegraphics[width=\linewidth]{figure/kvcache.pdf}  
  \caption{KVCache structure diagram}  
  \label{fig:kvcache}          
\end{figure}

\section{Background}
\label{background}

\subsection{KV Cache in Autoregressive Decoding}

Decoder-only LLM inference consists of two phases: prefill and decode.
During prefill, the model processes the input prompt in parallel and materializes the key and value states of all prompt tokens as the initial KV cache.
During decode, the model generates one token at a time; each new token appends its own key and value states to the KV cache of every layer, and subsequent tokens attend to the accumulated historical states.
\autoref{fig:kvcache} illustrates this process.
Unlike model weights, which are fixed once loaded, KV cache is dynamically created by active requests and grows with both prompt length and generated length.
For a batch of $B$ sequences with length $L$, its footprint scales as
$M_{\mathrm{KV}} \approx 2BLN_{\mathrm{layer}}H_{\mathrm{kv}}D_{\mathrm{head}}s$,
where the factor $2$ corresponds to keys and values, and $s$ is the number of bytes per element.

\subsection{Block-based KV Management in Modern LLM Servers}

Modern LLM serving systems commonly manage KV cache using a block-based memory abstraction.
Instead of allocating one large contiguous KV region for each request, the runtime partitions each request's logical KV sequence into fixed-size logical blocks and maps them to physical blocks from a global KV block pool.
A block table records this mapping, allowing a request's logical KV sequence to remain contiguous while its physical storage can be non-contiguous.
As shown in \autoref{fig:block}, when a request grows during decoding, the runtime appends newly allocated physical blocks to its block table; when a request finishes or its KV cache is reclaimed, the corresponding physical blocks are returned to the global pool.
This abstraction avoids large contiguous allocations, reduces fragmentation, and supports dynamic allocation, reclamation, and continuous batching for requests with different prompt and output lengths.

\begin{figure}[htbp]
\centering
\includegraphics[width=\linewidth]{figure/block.pdf}
\caption{Block-based KV management.}
\label{fig:block}
\end{figure}

\subsection{KV Sparsity in Attention}

In autoregressive decoding, each new query token attends to the historical KV cache.
Although dense attention considers all historical KV entries, their contributions are often highly uneven: only a small fraction of entries receive large attention weights or contribute significantly to the output representation, while many others have little influence.
This property is commonly referred to as KV sparsity.
KV sparsity is query-dependent and dynamic, because different generated tokens produce different queries and may rely on different historical entries.
It also varies across layers and attention heads, which capture different dependency patterns.
Therefore, KV importance should be viewed as a runtime property of attention computation rather than a static property of the input sequence.

\section{Bottleneck and Motivation}

\subsection{High-utilization Scheduling under Uncertain Output Length}

Modern LLM serving runtimes operate under a fixed GPU KV memory budget.
To improve batching efficiency and GPU throughput, the scheduler tends to keep this budget highly utilized by admitting requests and continuing running batches as long as the current KV footprint appears feasible.
In block-based KV management, this decision is made based on the current number of allocated KV blocks, the remaining free blocks, and a reserved safety margin.
However, such decisions are inherently fragile because output length is unknown at scheduling time.

When a request enters the system, its prompt length is known, but the number of tokens it will generate depends on the task, stopping criteria, sampling behavior, and user configuration.
Therefore, the scheduler can only reason about the current KV footprint and the immediate block demand of the next decode iteration.
A batch that fits the current budget may continue decoding for many more steps, gradually consuming the remaining KV block pool.
Once the free KV space is no longer sufficient to support the next decode iteration of the current batch, the runtime must preempt some active requests to recover KV blocks.

\autoref{fig:block-preempt} illustrates this behavior.
During decoding, KV usage repeatedly rises as active requests append new KV blocks and drops when preemption releases space.
The preemption rate also increases with both batch size and output length, showing that dynamic KV growth makes high-utilization scheduling increasingly unstable under heavier workloads.
This behavior is not caused by inefficient contiguous allocation; block-based runtimes already address that issue.
Rather, the key limitation is that high-utilization scheduling lacks a future-access-aware mechanism to proactively reduce or relocate KV blocks before the next decode iteration becomes unschedulable.

\begin{figure}[t]
\centering

\begin{subfigure}[t]{0.537\columnwidth}
\centering
\includegraphics[width=\linewidth]{figure/kv_usage_timeline.pdf}
\caption{KV usage timeline}
\label{fig:block-usage}
\end{subfigure}
\hfill
\begin{subfigure}[t]{0.443\columnwidth}
\centering
\includegraphics[width=\linewidth]{figure/preemptions_per_100_requests.pdf}
\caption{Request preemptions}
\label{fig:preempt}
\end{subfigure}

\caption{
Dynamic KV growth causes request preemption under high-utilization scheduling.
KV usage repeatedly increases during decoding until the remaining block budget cannot support the next decode iteration.
Preemptions become more frequent with larger batches and longer outputs.
}
\label{fig:block-preempt}
\end{figure}

\subsection{From Attention Sparsity to KV Lifecycle Management}

The opportunity for proactive KV management comes from the inherent sparsity of attention.
Under dense attention, every decode step accesses the full historical KV cache, making future KV access equivalent to a full scan of all accumulated KV entries.
In this case, there is little room for fine-grained lifecycle decisions beyond coarse-grained swapping or recomputation.

Attention sparsity changes this view.
For each decode step, generation quality can often be preserved by attending to only a small subset of historical KV entries.
We define this ideal subset as the minimal KV set required by that step.
Let $S_i^\star$ denote the minimal set of historical KV entries needed at decode step $i$ to preserve generation quality.
Because each generated token produces a different query, $S_i^\star$ changes over time.
The future sequence ${S_i^\star}$ therefore forms an ideal sparse-access trace of KV cache.

If the runtime knew this trace, it could derive an oracle lifecycle policy:
blocks that will be needed soon should remain resident or be prefetched;
blocks that are temporarily unused can be offloaded;
and blocks that never appear again can be permanently evicted.
Thus, the lifecycle value of a KV block is determined not only by whether it is important for the current step, but also by how it appears in future sparse-access sets.

\autoref{fig:trace} illustrates such sparse-access behavior across decoding steps.
The retained KV subset changes over time, showing that KV importance is dynamic rather than fixed after prefill.
However, the ideal future trace is unavailable at runtime because the server does not know future queries or the exact minimal KV set required by each future decode step.
This raises the central question: can practical sparse scoring provide useful evidence about this unavailable future sparse-access trace?

\begin{figure}[htbp]
\centering
\includegraphics[width=\linewidth]{figure/kv_trace.pdf}
\caption{
Sparse KV retention changes across decoding steps.
Each row shows retained prompt KV tokens under 50% compression on a LongBench sample.
The changing pattern indicates that KV importance is dynamic rather than fixed after prefill.
}
\label{fig:trace}
\end{figure}

\subsection{Can Practical Sparse Methods Approximate Future Ideal Sparse Access?}

We next examine whether practical sparse scoring can approximate the future ideal sparse-access trace.
The goal is not to predict arbitrary memory references, but to approximate the future minimal KV sets induced by attention sparsity.
Our measurements reveal two properties that make this approximation feasible for online KV lifecycle management.

\noindent\textbf{Observation 1: Sparse KV sets are locally stable but not static.}
\autoref{fig:kv-set-overlap} shows the overlap between important KV sets at step $i$ and step $i+\Delta$.
The overlap is high when $\Delta$ is small, indicating that nearby decode steps often require similar KV entries.
This local stability allows current sparse scores to serve as useful predictors of near-future KV importance.
For example, blocks that are important in the current step are more likely to remain useful in the next few steps, making it reasonable to keep them resident or prefetch them when they are offloaded.

At the same time, the overlap decreases as $\Delta$ grows.
This indicates that KV importance is not static throughout generation.
A one-shot compression decision cannot reliably represent sparse access over a long horizon.
Instead, sparse scores must be refreshed online so that the runtime can continuously track the evolving KV importance pattern.

\begin{figure}[htbp]
\centering
\includegraphics[width=\linewidth]{figure/kv_overlap.pdf}
\caption{
Sparse KV sets are locally stable but not static.
Overlap is high for nearby decoding steps and decreases as step distance grows, motivating online score refresh.
}
\label{fig:kv-set-overlap}
\end{figure}

\noindent\textbf{Observation 2: Sparse eviction is spatially continuous.}
\autoref{fig:eviction} visualizes token-level sparse eviction decisions over long-context inputs.
Evicted KV entries are not uniformly scattered across the sequence; instead, they often form large contiguous spans.
This spatial continuity suggests that neighboring tokens tend to have correlated KV importance.
It is especially important for LLM serving systems because modern runtimes allocate, reclaim, offload, and prefetch KV cache at block granularity.
Therefore, token-level sparse evidence can be lifted to block-level scoring without losing the dominant sparse structure, while directly matching the runtime's management unit.

\begin{figure}[hbpt]
\centering

\begin{subfigure}[t]{0.515\columnwidth}
\centering
\includegraphics[width=\linewidth]{figure/eviction_mask_heatmap.pdf}
\caption{Eviction mask}
\label{fig:eviction-heatmap}
\end{subfigure}
\hfill
\begin{subfigure}[t]{0.468\columnwidth}
\centering
\includegraphics[width=\linewidth]{figure/evicted_run_length_vs_random.pdf}
\caption{Evicted span length}
\label{fig:evicted-length}
\end{subfigure}

\caption{
Sparse eviction is spatially clustered.
Token-level eviction decisions form contiguous spans, allowing sparse evidence to be lifted to block granularity.
}
\label{fig:eviction}
\end{figure}

\subsection{Opportunity: Sparse Scores as KV Lifecycle Hints}

The observations above suggest that practical sparse scoring can serve as an online approximation of the future ideal sparse-access trace.
Local stability makes current sparse scores predictive of near-future KV value, while long-horizon drift shows that this approximation must be refreshed during decoding.
Spatial continuity further makes the approximation compatible with block-based KV managers.
Together, these properties turn sparse scoring from a local attention optimization into a runtime-level signal for KV lifecycle management.
Instead of only deciding which KV entries participate in the current attention computation, sparse scores can guide whether KV blocks should remain on GPU, be prefetched, be temporarily offloaded, or be permanently evicted.
This motivates KVCore, a block-level KV lifecycle management system that uses practical sparse scoring to approximate future ideal sparse access and proactively reduce GPU KV pressure.


\subsection{Challenges}

The motivation above suggests that practical sparse scores can approximate future sparse access.
Turning this signal into a serving-system mechanism requires addressing two challenges:
deriving block-level, low-overhead importance signals, and making safe lifecycle decisions under imperfect prediction, memory pressure, and data-movement overheads.

\subsubsection{Block-level Sparse Approximation}

The first challenge is deriving a runtime-friendly KV importance signal.
The ideal sparse-access trace ${S_i^\star}$ is unavailable at runtime, and practical sparse methods must approximate it online.
However, many existing methods output token- or chunk-level sparse decisions, while modern serving systems allocate, reclaim, offload, prefetch, and evict KV cache at block granularity.
A lifecycle manager therefore needs block-level scores that preserve useful fine-grained sparse evidence while directly matching the runtime's management unit.

The scoring path must also be lightweight and backend-independent.
Methods that materialize full attention scores introduce high overhead and are tightly coupled to specific attention kernels.
Moreover, lifecycle actions such as prefetching and offloading require importance signals early enough to guide data movement before they block future attention computation.
Thus, the challenge is to convert sparse-access approximation into a query-aware, continuously refreshed, and low-overhead block score that can be directly consumed by a production-style KV manager.

\subsubsection{Lifecycle Decisions under Imperfect Prediction}

The second challenge is using approximate sparse scores without harming quality or throughput.
A low-scored block may be permanently useless, or it may simply be temporarily cold and become important again soon.
Similarly, a high-scored block may be needed immediately but not reused over a longer horizon.
Thus, sparse scores cannot be treated as oracle keep/drop decisions.

Different lifecycle actions expose different trade-offs.
Permanent eviction saves the most GPU memory but risks quality loss under wrong prediction.
Temporary offload is recoverable but introduces CPU--GPU transfer overhead.
Prefetching can hide future transfer latency but wastes bandwidth and GPU blocks if inaccurate.
Keeping blocks resident is safest for quality but increases GPU KV pressure.
A robust runtime must therefore combine sparse scores with score history and memory pressure, choose conservative actions when confidence is low, and overlap block movement with computation to avoid replacing preemption overhead with transfer stalls.



\section{Design}
\label{sec:design}

\subsection{Overview}

KVCore is a standalone block-based LLM serving framework that integrates sparse scoring and KV lifecycle management as first-class runtime components.
As shown in \autoref{fig:kvcore-overview}, KVCore follows a vLLM-style execution model with an engine core, scheduler, KV manager, model runner, and attention backend.
The scheduler forms execution batches from waiting and running requests.
The KV manager maintains per-request and per-layer logical block tables, allocates physical blocks from a global GPU KV block pool, and tracks each block's lifecycle state.
The model runner builds input batches, slot mappings, sparse read block tables, and paged-attention metadata, while the attention backend writes newly generated K/V states and reads historical KV blocks through the provided block table.

\begin{figure}[htbp]
\centering
\includegraphics[width=\linewidth]{figure/arch.pdf}
\caption{KVCore system overview.}
\label{fig:kvcore-overview}
\end{figure}

KVCore separates sparse attention planning from lifecycle transitions.
In the \emph{data path}, KVCore selects the blocks visible to the current attention computation and materializes a sparse read block table.
This path changes only per-forward metadata and does not necessarily release physical KV blocks.
In the \emph{lifecycle path}, KVCore uses sparse scores, score history, request state, and global memory pressure to decide whether a block should remain on GPU, become compute-cold, be offloaded to CPU, be prefetched back to GPU, or be permanently evicted.
This separation prevents KVCore from treating every low score as an irreversible deletion.
At scheduler boundaries, KVCore commits block-state updates and block-table mutations, keeping them away from in-flight attention computation while allowing scoring and data movement to overlap with serving execution.

\subsection{Block-level Sparse Scoring}

\begin{algorithm}[t]
\caption{\textsc{KVCore}: Low-Cost Query-Guided Block Compression}
\label{alg:blockwise-press}
\small
\begin{algorithmic}[1]
\Require Key/value cache $K,V \in \mathbb{R}^{H \times T \times d}$ at one layer,
hidden states $X$, block size $B$, compression ratio $\rho$,
query window $w$, protected prefix blocks $p$, protected recent blocks $r$.
\Ensure Compressed cache $\widetilde{K},\widetilde{V}$.
\State Partition $K,V$ into $n=\lceil T/B\rceil$ contiguous blocks
       $\{(K_i,V_i)\}_{i=1}^{n}$.
\State Obtain the last-$w$ pre-RoPE query states
       $Q \leftarrow \mathrm{Proj}_{Q}(X_{T-w:T})$.
\For{$i=1$ \textbf{to} $n$}
    \State Compute mean anchor
       $\mu_i \leftarrow \frac{1}{|K_i|}\sum_{k \in K_i} k$.
    \State Select $m$ high-norm keys
       $R_i \leftarrow \mathrm{TopM}_{k\in K_i} \|k\|_2$,
       and compute top-key anchor $\tau_i \leftarrow \frac{1}{m}\sum_{r \in R_i} r$.
    \State Score the mean anchor:
       $a_{h,i} \leftarrow
       \max_{q \in Q}
       \frac{\langle q_h,\mu_{i,h}\rangle}{\sqrt{d}}$.
    \State Score the top-key anchor:
       $c_{h,i} \leftarrow
       \max_{q \in Q}
       \frac{\langle q_h,\tau_{i,h}\rangle}{\sqrt{d}}$.
    \State Fuse the two anchor scores:
       $s_{h,i} \leftarrow \lambda a_{h,i} + (1-\lambda)c_{h,i}$.
    \State Aggregate heads:
       $s_i \leftarrow \frac{1}{H}\sum_{h=1}^{H}s_{h,i}$.
\EndFor
\State Set block budget $b \leftarrow \lceil n(1-\rho)\rceil$.
\State Initialize protected set
       $\mathcal{P} \leftarrow \{1,\ldots,p\}
       \cup \{n-r+1,\ldots,n\}$.
\If{the last block is partial}
    \State $\mathcal{P} \leftarrow \mathcal{P}\cup\{n\}$.
\EndIf
\If{$|\mathcal{P}| \le b$}
    \State Select $\mathcal{S} \leftarrow \mathcal{P}
    \cup \mathrm{Top}_{b-|\mathcal{P}|}
    \left(\{s_i: i \notin \mathcal{P}\}\right)$.
\Else
    \State Select $\mathcal{S} \leftarrow
    \mathrm{Top}_{b}\left(\{s_i: 1\le i\le n\}\right)$.
\EndIf
\State Expand selected blocks to token indices
       $\mathcal{I}\leftarrow \bigcup_{i\in\mathcal{S}}\{(i-1)B,\ldots,\min(iB,T)-1\}$.
\State \Return $\widetilde{K}\leftarrow K[:,\mathcal{I},:]$,
       $\widetilde{V}\leftarrow V[:,\mathcal{I},:]$.
\end{algorithmic}
\end{algorithm}

\begin{table}[t]
\centering
\small
\caption{KVCore main configuration.}
\label{tab:blockwise-config}
\begin{tabular}{lc}
\toprule
Parameter & Value \\
\midrule
Block size $B$ & 16 \\
Query window $w$ & 64 \\
Summary & Mean+TopK \\
Top-key selector & KeyNorm \\
Top keys per block $m$ & 4 \\
Mean-key weight $\lambda$ & 0.75 \\
Query aggregation & Max \\
Head aggregation & Mean \\
Protected prefix blocks & 1 \\
Protected recent blocks & 2 \\
\bottomrule
\end{tabular}
\end{table}

KVCore needs a sparse signal that is accurate enough to approximate important KV blocks and cheap enough to run inside an online serving framework.
Algorithm~\ref{alg:blockwise-press} gives the scoring procedure, and \autoref{tab:blockwise-config} lists the default configuration.
The method follows the intuition used by prior KV compression methods such as SnapKV~\cite{snapkv}: recent queries provide evidence for identifying important historical KV entries.
KVCore adapts this idea to block-based serving by computing query-guided scores over lightweight block summaries instead of materializing full token-level attention scores.

\noindent\textbf{Query-guided block scoring.}
For each completed KV block, KVCore builds two key-side anchors: a mean anchor that captures stable block-level relevance, and a top-key anchor that averages $m$ high-norm keys to preserve high-response representatives.
These summaries are compact, reusable across refresh intervals, and independent of attention-kernel internals.
Incomplete tail blocks are protected until their contents become stable.

KVCore scores blocks using the most recent query window.
Given the last $w$ pre-RoPE query states, it aligns query heads with KV heads for GQA/MQA models and computes query-summary dot products for the two anchors.
It then uses max aggregation over the query window to capture cases where any recent query strongly depends on a block.
The mean-anchor and top-key-anchor scores are fused by weight $\lambda$, and head scores are averaged into one block score.
The resulting score is a block-level approximation of sparse access value, not an oracle prediction of future attention.

\noindent\textbf{Protected blocks.}
KVCore protects structurally important blocks before selecting top-scored blocks.
Following observations on attention sinks and recent-token locality~\cite{xiao2024efficient,zhang2023h2o,li2024snapkv}, it protects a small number of prefix blocks and recent blocks.
The current tail block and filling blocks are also protected for correctness, because they may not have stable summaries or complete physical contents.
These protected regions reduce the risk that sparse selection removes global anchors or local context.
Algorithm~\ref{alg} presents the block selection procedure with a fixed compression ratio $\rho$ for controlled evaluation, while the serving system determines the actual active budget adaptively as part of lifecycle management.

\noindent\textbf{Score history and refresh.}
Sparse scores depend on the current query window and may fluctuate across steps.
KVCore therefore maintains an exponential moving average for each logical block to smooth query-level noise and avoid high-risk actions from a single low score.
Scores are refreshed periodically: the interval is shortened under high pressure or score divergence, and lengthened under low pressure to reduce overhead.

\begin{figure}[hbpt]
\centering

\begin{subfigure}[t]{0.49\columnwidth}
\centering
\includegraphics[width=\linewidth]{figure/kv_head_score_cosine.pdf}
\caption{Head similarity}
\label{fig:head-sim}
\end{subfigure}
\hfill
\begin{subfigure}[t]{0.49\columnwidth}
\centering
\includegraphics[width=\linewidth]{figure/layer_block_index_jaccard.pdf}
\caption{Layer index similarity}
\label{fig:layer-index-sim}
\end{subfigure}

\caption{
Similarity of sparse scoring structures across heads and layers.
(a) Cosine similarity of block-summary scores across KV head groups.
(b) Jaccard similarity of block indices selected by KVCore across layers.
High similarity indicates that sparse indices can be reused across similar heads and layer regions to reduce scoring overhead.
}
\label{fig:reuse}
\end{figure}

\noindent\textbf{Sharing scores across heads and layers.}
KVCore further reduces scoring overhead by exploiting structural similarity across heads and layers.
As shown in \autoref{fig:head-sim}, many heads produce similar block-score patterns, so KVCore groups similar heads and shares summaries or block scores within each group.
\autoref{fig:layer-index-sim} further shows that layer-wise selected block indices are not uniformly similar across the whole model; instead, they form several clear high-similarity regions separated by lower-similarity boundaries.
KVCore uses these boundaries to partition layers into reuse regions and selects one anchor layer per region.
The block indices selected for the anchor layer are reused by other layers in the same region, reducing per-layer scoring overhead while avoiding reuse across dissimilar layer regions.
This sharing is only an optimization: KVCore periodically recalibrates groups with full scoring and disables reuse when divergence becomes large.

\subsection{Sparsity-guided KV Lifecycle Management}

The lifecycle manager turns block-level sparse scores into runtime memory actions.
KVCore deliberately separates two decisions.
The first is \emph{dynamic sparse computation}, which decides which blocks are visible to the next attention computation.
This decision only changes the sparse read block table and is reversible.
The second is \emph{permanent reclamation}, which decides which blocks can be irreversibly removed when the system must release GPU KV blocks.
This separation keeps the common sparse-compute path simple while making permanent eviction conservative.

\noindent\textbf{Lifecycle states.}
KVCore tracks each logical KV block using six states, as shown in \autoref{tab:lifecycle-states}.
A block can be active on GPU, cold on GPU, stored on CPU, in transfer, or permanently evicted.
There is no separate warm state: once a prefetched block finishes its CPU-to-GPU transfer, it directly becomes GPU-active and can be referenced by the sparse read block table.

\begin{table}[t]
\centering
\small
\caption{KV block lifecycle states in KVCore.}
\label{tab:lifecycle-states}
\begin{tabular}{lll}
\toprule
State & Location & Role \\
\midrule
GPU-active & GPU & selected by sparse read table \\
GPU-cold & GPU & resident but not selected \\
CPU & CPU & offloaded, recoverable \\
Offloading & GPU$\rightarrow$CPU & async offload pending \\
Prefetching & CPU$\rightarrow$GPU & async prefetch pending \\
Evicted & none & null block, unrecoverable \\
\bottomrule
\end{tabular}
\end{table}

The main state transitions are simple.
After sparse planning, selected GPU-resident blocks become GPU-active, while unselected GPU-resident blocks become GPU-cold.
GPU-cold blocks can be asynchronously offloaded to CPU or permanently evicted.
When a CPU block becomes hot again, KVCore marks it as Prefetching; after the transfer completes, it becomes GPU-active.
Evicted blocks are replaced by null entries in the logical block table and are excluded from future sparse plans.

\noindent\textbf{Block-completion score refresh.}
KVCore refreshes sparse scores only at block-completion boundaries.
When a logical KV block becomes full, KVCore builds its block summary, evaluates block scores using the recent query window, and updates the history score of each valid block.
Between two block-completion boundaries, the current sparse plan is reused and the filling tail block is protected.
This aligns scoring with the runtime's block management unit and avoids per-token lifecycle decisions.

For each block, KVCore maintains a current score and a history score.
The current score reflects the block's importance for the latest query window.
The history score is a running average that captures long-term usefulness.
KVCore also records the inactive age of each block, defined as the number of score-refresh intervals since the block was last selected as GPU-active.
These three signals are sufficient for lifecycle decisions:
current score identifies blocks that are hot now, while history score and inactive age identify blocks that have remained cold for a long time.

\noindent\textbf{Dynamic sparse computation.}
Dynamic sparse computation uses only the current block scores.
For each request and layer, KVCore first normalizes valid block scores into non-negative weights.
It then sorts blocks by current score and selects the smallest set whose cumulative normalized score reaches a fixed top-p threshold.
By default, KVCore uses $p=0.9$:

$$
\sum_{b \in \mathcal{A}*{r,l}} \hat{s}*{r,l,b}^{t} \ge p,
\quad p=0.9 .
$$

Here, $\mathcal{A}*{r,l}$ is the active set for request $r$ and layer $l$, and $\hat{s}*{r,l,b}^{t}$ is the normalized current score of block $b$.
Protected prefix, recent, and filling blocks are always included in the active set.
The selected blocks are inserted into the sparse read block table.
If a selected block is already on GPU, it becomes GPU-active.
If it is on CPU, KVCore marks it as Prefetching and schedules a CPU-to-GPU transfer.
GPU-resident blocks that are not selected become GPU-cold.

This top-p rule avoids a fixed compression ratio.
When scores are concentrated, a small number of blocks can cover most of the importance mass.
When scores are more diffuse, KVCore naturally keeps more blocks.
Because dynamic sparse computation is reversible, KVCore can update the active set whenever scores are refreshed without permanently deleting unselected blocks.

\noindent\textbf{Fair permanent eviction.}
Permanent eviction is triggered only when the scheduler requires additional free GPU blocks.
KVCore reuses the scheduler's free-block target and computes the number of blocks that must be released:

$$
M_{\mathrm{release}}
=
\max\left(0, F_{\mathrm{target}} - F_t\right).
$$

where $F_t$ is the current number of free GPU blocks and $F_{\mathrm{target}}$ is the target number of free blocks required by the scheduler for safe execution.
If $M_{\mathrm{release}}=0$, KVCore does not permanently evict blocks.

When permanent eviction is necessary, KVCore distributes the eviction quota across requests according to their existing eviction ratios.
For request $r$, let $E_r$ be the number of blocks that have already been permanently evicted, and let $C_r$ be the number of completed valid blocks.
KVCore defines the request-level eviction ratio as

$$
e_r = \frac{E_r}{C_r}.
$$

Requests with lower eviction ratios receive higher priority for contributing eviction candidates, while requests that have already been heavily evicted are protected from further irreversible loss.
Let $C_r^{\mathrm{cold}}$ be the number of evictable GPU-cold blocks of request $r$.
KVCore assigns request $r$ an eviction quota

$$
q_r =
\left\lceil
M_{\mathrm{release}}
\cdot
\frac{(1-e_r) C_r^{\mathrm{cold}}}
{\sum_{r'} (1-e_{r'}) C_{r'}^{\mathrm{cold}}}
\right\rceil .
$$

Within each request, KVCore ranks evictable GPU-cold blocks using a simple lexicographic rule:
lower history score first, longer inactive age first, and lower current score first.
This selects blocks that have remained unimportant for a long time.
Protected blocks and filling blocks are never evicted.
Compared with using only the current score, this rule avoids deleting blocks that are temporarily cold but historically useful.

\begin{algorithm}[t]
\caption{\textsc{KVCore}: Sparse-guided KV Lifecycle Update}
\label{alg:lifecycle}
\small
\begin{algorithmic}[1]
\Require Running requests $\mathcal{R}$, block states, current scores $s$, history scores $\bar{s}$, free blocks $F_t$, scheduler target $F_{\mathrm{target}}$, top-p threshold $p=0.9$.
\Ensure Sparse read block tables and lifecycle actions.
\State Compute required release $M_{\mathrm{release}} \leftarrow \max(0, F_{\mathrm{target}} - F_t)$.

\ForAll{request $r \in \mathcal{R}$ and layer $l$}
\If{a logical block becomes full}
\State Build the new block summary.
\State Refresh current scores $s_{r,l,b}$ using the recent query window.
\State Update history scores $\bar{s}_{r,l,b}$ and inactive ages.
\EndIf
\State Select active set $\mathcal{A}_{r,l}$ by top-p cumulative current score.
\State Force protected prefix, recent, and filling blocks into $\mathcal{A}_{r,l}$.
\State Build the sparse read block table from $\mathcal{A}*{r,l}$.
\State Mark selected GPU blocks as GPU-active.
\State Mark selected CPU blocks as Prefetching.
\State Mark unselected GPU-resident blocks as GPU-cold.
\EndFor

\If{$M_{\mathrm{release}} > 0$}
\State Compute request eviction ratios $e_r$ and quotas $q_r$.
\ForAll{request $r \in \mathcal{R}$}
\State Rank evictable GPU-cold blocks by lower history score, longer inactive age, and lower current score.
\State Permanently evict up to $q_r$ ranked blocks.
\EndFor
\EndIf

\State Offload remaining GPU-cold blocks that are not evicted.
\State Prefetch CPU blocks marked as Prefetching.
\State Commit block-state and block-table updates at scheduler boundaries.
\end{algorithmic}
\end{algorithm}

\noindent\textbf{Concurrent offload and prefetch.}
Permanent eviction only handles blocks that must be irreversibly removed to satisfy the scheduler's release demand.
All other cold and hot transitions are handled by asynchronous movement.
After each score refresh, GPU blocks that become cold are offloaded to CPU unless they are selected for permanent eviction.
At the same time, CPU blocks that enter the new active set are prefetched back to GPU.
Thus, offload and prefetch can proceed concurrently: the offload path releases GPU capacity, while the prefetch path restores blocks that are predicted to be useful.

KVCore uses separate CUDA streams for offload and prefetch.
Offload moves GPU-cold blocks through the Offloading state and returns their GPU physical blocks to the free pool only after the GPU-to-CPU copy completes.
Prefetch moves CPU blocks through the Prefetching state and changes them directly to GPU-active after the CPU-to-GPU copy completes.
A sparse read block table can reference a prefetched block only after its transfer event has completed; otherwise, the compute stream waits on the event or the block is delayed to the next sparse plan.
All block-table mutations are committed at scheduler boundaries, so in-flight attention computation never observes partially updated block states.

Overall, KVCore uses current scores for reversible sparse computation, historical scores and inactive age for conservative permanent eviction, and asynchronous offload/prefetch for recoverable movement.
This keeps the lifecycle policy simple while still allowing the runtime to reduce GPU KV pressure before request preemption occurs.




\section{Evaluation}
\subsection{Experimental Setup}

\textbf{System Configuration.} 
All experiments were conducted on a single-machine setup running Ubuntu 20.04. 
The computing node is equipped with an Intel(R) Xeon(R) W-3345 CPU @ 3.00GHz, 128GB DDR4 memory, and an NVIDIA L40s GPU with 48GB of VRAM, connected via PCIe Gen4 x16. 

\begin{table}[h!]
\centering
\caption{System configuration for the experiments.}
\begin{tabular}{ll}
\hline
Component & Specification \\
\hline
OS & Ubuntu 20.04 (kernel 5.15.0) \\
CPU & Intel(R) Xeon(R) W-3345 @ 3.00GHz \\
Memory & 128GB DDR4 \\
GPU & NVIDIA L40s, 48GB VRAM, PCIe Gen4 x16 \\
\hline
\end{tabular}
\end{table}

\textbf{Models.} 
The experiments used three representative large language models: LLaMA-3.1-8B-Instruct, Qwen3-8B, and Mistral-7B-Instruct-v0.3. 
These models were selected to evaluate the performance of the proposed compression algorithms across different architectures and system characteristics.

\textbf{Datasets and Metrics.} 
The experimental datasets include all 16 subtasks of LongBench, used to evaluate model accuracy on long-context reasoning tasks; PG19, used for perplexity evaluation on LLaMA-3.1-8B-Instruct. 
Evaluation metrics include accuracy, perplexity, end-to-end system performance (throughput and latency), and preemption rate, providing a comprehensive assessment of both model reasoning capability and system scheduling efficiency.

\textbf{Baselines.} 
At the algorithm level, the proposed method was compared with two KV compression approaches: SnapKV, a classical token-wise KV compression method representing traditional fine-grained sparse strategies, and ChunkKV, a state-of-the-art block-wise KV compression method representing the latest development in block-level compression. 
At the system level, comparisons were made with vLLM and InfiniGen. 
vLLM is a high-performance LLM inference system supporting long-context reasoning, employing block-wise KV management and asynchronous scheduling; InfiniGen is a lightweight framework designed for high-throughput inference, capable of handling large-scale continuous requests. 
These baselines allow evaluation of the proposed method from both algorithmic and system perspectives.

\subsection{Overall Performance}

\begin{table*}[!htbp]
\centering
\scriptsize
\setlength{\tabcolsep}{2.0pt}
\renewcommand{\arraystretch}{1.08}
\resizebox{\textwidth}{!}{%
\begin{tabular}{lllrrrrrrrrrrrrrrrrr}
\toprule
Model & Ratio & Method & \multicolumn{3}{c}{Single-Doc QA} & \multicolumn{3}{c}{Multi-Doc QA} & \multicolumn{3}{c}{Summarization} & \multicolumn{3}{c}{Few-shot} & \multicolumn{2}{c}{Synthetic} & \multicolumn{2}{c}{Code} & Avg. \\
\cmidrule(lr){4-6}\cmidrule(lr){7-9}\cmidrule(lr){10-12}\cmidrule(lr){13-15}\cmidrule(lr){16-17}\cmidrule(lr){18-19}
 &  &  & \rotatebox{60}{NarrativeQA} & \rotatebox{60}{Qasper} & \rotatebox{60}{MultiFieldQA} & \rotatebox{60}{HotpotQA} & \rotatebox{60}{2WikiMQA} & \rotatebox{60}{MuSiQue} & \rotatebox{60}{GovReport} & \rotatebox{60}{QMSum} & \rotatebox{60}{MultiNews} & \rotatebox{60}{TriviaQA} & \rotatebox{60}{SAMSum} & \rotatebox{60}{TREC} & \rotatebox{60}{PassageCnt} & \rotatebox{60}{PassageRet} & \rotatebox{60}{LCC} & \rotatebox{60}{RepoBench-P} & Avg. \\
\midrule
\multicolumn{20}{c}{\textbf{Llama-3.1-8B-Instruct}} \\
\midrule
 & 0.3 & KVCore & 30.07 & \textbf{48.09} & \textbf{56.20} & 58.82 & 50.04 & 32.69 & 33.99 & 24.66 & 26.29 & \textbf{92.23} & 39.59 & 28.50 & 11.05 & \textbf{100.00} & 53.04 & \textbf{49.23} & 45.91 \\
 &  & SnapKV & 30.69 & 46.77 & 55.48 & \textbf{59.43} & \textbf{51.10} & 32.39 & 33.71 & 24.75 & \textbf{26.35} & 91.41 & \textbf{41.04} & \textbf{31.50} & \textbf{11.20} & \textbf{100.00} & \textbf{53.87} & 47.52 & \textbf{46.08} \\
 &  & ChunkKV & \textbf{30.76} & 46.83 & 55.65 & 58.31 & 50.10 & \textbf{33.93} & \textbf{34.19} & \textbf{24.92} & 26.31 & 91.41 & 40.47 & 28.50 & 11.15 & \textbf{100.00} & 52.46 & 47.94 & 45.81 \\
\addlinespace[1pt]
 & 0.4 & KVCore & 29.58 & 47.87 & \textbf{55.86} & 58.96 & 49.05 & 31.67 & \textbf{33.71} & 24.56 & 26.08 & 91.54 & 38.87 & 24.50 & 11.13 & \textbf{100.00} & 52.92 & \textbf{48.39} & 45.29 \\
 &  & SnapKV & 30.21 & 47.11 & 54.83 & \textbf{59.19} & \textbf{51.71} & \textbf{33.31} & 33.50 & 24.66 & 26.12 & \textbf{91.88} & \textbf{40.69} & \textbf{37.50} & 10.65 & \textbf{100.00} & \textbf{53.04} & 47.25 & \textbf{46.35} \\
 &  & ChunkKV & \textbf{30.61} & \textbf{48.08} & 55.72 & 59.16 & 51.54 & 32.77 & 33.62 & \textbf{24.79} & \textbf{26.53} & 91.41 & 40.46 & 24.50 & \textbf{12.55} & \textbf{100.00} & 52.44 & 48.33 & 45.78 \\
\addlinespace[1pt]
 & 0.5 & KVCore & 29.38 & 47.34 & 54.59 & 58.23 & 49.52 & 31.30 & 32.83 & 24.47 & 25.55 & 91.72 & 39.43 & 23.00 & 11.05 & \textbf{100.00} & 52.02 & \textbf{48.93} & 44.96 \\
 &  & SnapKV & \textbf{30.79} & 46.22 & \textbf{55.49} & 59.14 & \textbf{50.98} & \textbf{32.84} & 32.64 & 24.51 & \textbf{25.66} & \textbf{91.88} & \textbf{40.56} & \textbf{36.50} & \textbf{11.65} & \textbf{100.00} & \textbf{53.09} & 47.17 & \textbf{46.20} \\
 &  & ChunkKV & 30.44 & \textbf{47.46} & 54.68 & \textbf{59.16} & 49.77 & 32.38 & \textbf{33.43} & \textbf{24.81} & 25.64 & 91.43 & 39.67 & 24.00 & 10.55 & \textbf{100.00} & 50.81 & 48.54 & 45.17 \\
\addlinespace[1pt]
 & 0.6 & KVCore & 27.38 & \textbf{47.97} & 51.83 & 57.51 & 49.62 & 31.56 & 31.81 & 23.85 & \textbf{24.73} & \textbf{91.98} & 38.00 & 24.00 & 7.50 & 99.50 & 51.22 & \textbf{50.27} & 44.30 \\
 &  & SnapKV & \textbf{30.60} & 45.46 & \textbf{55.14} & \textbf{59.25} & \textbf{51.15} & \textbf{32.94} & 31.41 & \textbf{24.86} & 24.64 & 91.54 & 39.85 & \textbf{36.50} & 11.70 & \textbf{100.00} & \textbf{53.28} & 47.15 & \textbf{45.97} \\
 &  & ChunkKV & 30.20 & 47.08 & 54.96 & 58.33 & 48.70 & 32.87 & \textbf{32.35} & 24.39 & 24.52 & \textbf{91.98} & \textbf{40.39} & 26.50 & \textbf{12.05} & \textbf{100.00} & 51.18 & 48.86 & 45.27 \\
\addlinespace[1pt]
 & 0.7 & KVCore & 26.99 & \textbf{46.53} & 50.43 & 57.97 & 49.02 & 31.01 & 30.56 & 23.96 & \textbf{24.35} & 91.48 & 36.60 & 22.00 & 6.54 & 99.50 & 51.08 & \textbf{50.48} & 43.66 \\
 &  & SnapKV & \textbf{30.03} & 46.04 & \textbf{55.84} & \textbf{59.04} & \textbf{50.92} & \textbf{32.56} & 30.47 & \textbf{24.49} & 23.87 & \textbf{92.08} & \textbf{40.66} & \textbf{41.00} & 11.63 & \textbf{100.00} & \textbf{53.32} & 47.52 & \textbf{46.22} \\
 &  & ChunkKV & 29.63 & 46.35 & 55.66 & 58.87 & 47.79 & 32.39 & \textbf{31.08} & 24.34 & 23.46 & 91.68 & 38.74 & 24.00 & \textbf{12.55} & 99.50 & 50.83 & 49.13 & 44.75 \\
\midrule
\multicolumn{20}{c}{\textbf{Mistral-7B-Instruct-v0.3}} \\
\midrule
 & 0.3 & KVCore & 26.95 & 38.57 & 51.33 & \textbf{48.97} & \textbf{39.46} & 27.84 & 33.82 & 25.34 & \textbf{26.50} & 86.02 & \textbf{24.06} & \textbf{50.73} & 4.73 & 98.00 & \textbf{51.57} & 56.46 & \textbf{43.15} \\
 &  & SnapKV & 27.95 & 38.44 & 51.01 & 48.75 & 38.80 & 28.11 & 33.85 & \textbf{25.57} & 26.08 & \textbf{86.49} & 22.79 & 49.63 & 5.17 & 98.00 & 51.14 & \textbf{56.51} & 43.02 \\
 &  & ChunkKV & \textbf{28.07} & \textbf{39.41} & \textbf{51.70} & 48.60 & 38.40 & \textbf{28.93} & \textbf{33.99} & 25.29 & 26.01 & 85.42 & 22.58 & 48.92 & \textbf{6.12} & \textbf{98.50} & 51.54 & 56.03 & 43.09 \\
\addlinespace[1pt]
 & 0.4 & KVCore & 26.27 & 38.60 & 51.67 & 48.54 & \textbf{39.62} & 27.61 & 33.51 & 24.85 & 25.98 & \textbf{86.52} & \textbf{24.57} & \textbf{51.76} & 5.25 & 98.00 & \textbf{52.51} & 56.33 & \textbf{43.22} \\
 &  & SnapKV & 27.42 & 38.55 & 51.21 & 48.71 & 38.89 & \textbf{28.41} & 33.23 & \textbf{25.26} & \textbf{26.07} & 85.99 & 22.69 & 49.12 & \textbf{5.77} & \textbf{98.50} & 52.50 & \textbf{56.82} & 43.07 \\
 &  & ChunkKV & \textbf{27.54} & \textbf{39.61} & \textbf{51.70} & \textbf{48.90} & 37.18 & 28.34 & \textbf{33.63} & 24.86 & 25.83 & 85.96 & 22.80 & 48.10 & 4.92 & 98.00 & 51.44 & 56.47 & 42.83 \\
\addlinespace[1pt]
 & 0.5 & KVCore & 25.65 & 37.37 & 50.51 & 47.82 & 36.35 & 27.19 & \textbf{33.14} & 24.06 & 25.44 & \textbf{86.92} & \textbf{24.74} & \textbf{51.42} & 4.53 & \textbf{98.50} & 52.23 & 56.46 & 42.65 \\
 &  & SnapKV & 27.23 & 38.47 & \textbf{51.14} & 48.01 & \textbf{38.74} & \textbf{27.99} & 32.60 & \textbf{25.24} & \textbf{25.48} & 86.56 & 23.25 & 49.63 & 4.85 & \textbf{98.50} & 52.52 & \textbf{57.72} & \textbf{43.00} \\
 &  & ChunkKV & \textbf{27.52} & \textbf{39.87} & 50.67 & \textbf{48.34} & 37.41 & 27.95 & 32.95 & 24.94 & \textbf{25.48} & 86.23 & 24.23 & 45.67 & \textbf{5.27} & 98.00 & \textbf{52.78} & 57.24 & 42.78 \\
\addlinespace[1pt]
 & 0.6 & KVCore & 25.36 & \textbf{38.13} & 50.45 & \textbf{48.59} & 36.01 & 25.16 & \textbf{31.97} & 23.63 & \textbf{25.08} & \textbf{87.66} & \textbf{26.26} & \textbf{51.74} & 4.03 & \textbf{98.50} & 52.09 & 57.19 & 42.62 \\
 &  & SnapKV & 27.11 & 37.06 & \textbf{51.02} & 48.58 & \textbf{38.28} & 27.37 & 31.51 & \textbf{24.88} & 24.62 & 87.00 & 23.58 & 48.38 & 4.01 & \textbf{98.50} & \textbf{53.38} & \textbf{58.03} & \textbf{42.71} \\
 &  & ChunkKV & \textbf{27.63} & 37.68 & 50.40 & 47.42 & 36.38 & \textbf{27.83} & 31.69 & 24.63 & 24.65 & 86.26 & 25.91 & 40.89 & \textbf{5.26} & \textbf{98.50} & 52.88 & 57.40 & 42.21 \\
\addlinespace[1pt]
 & 0.7 & KVCore & 25.84 & 35.74 & 49.62 & \textbf{48.31} & 34.74 & 25.23 & 30.51 & 23.11 & \textbf{24.28} & \textbf{87.85} & \textbf{24.98} & \textbf{51.04} & 2.84 & 97.50 & 51.79 & 56.78 & 41.88 \\
 &  & SnapKV & 27.00 & 35.82 & \textbf{50.69} & 47.69 & \textbf{37.67} & 27.70 & 30.28 & \textbf{24.98} & 23.95 & 86.94 & 24.66 & 48.50 & 4.68 & \textbf{98.50} & 53.34 & \textbf{57.99} & \textbf{42.52} \\
 &  & ChunkKV & \textbf{27.58} & \textbf{36.39} & 50.20 & 47.48 & 35.69 & \textbf{28.00} & \textbf{30.89} & 24.38 & 23.47 & 86.76 & 24.95 & 36.75 & \textbf{5.55} & \textbf{98.50} & \textbf{54.07} & 57.26 & 41.74 \\
\midrule
\multicolumn{20}{c}{\textbf{Qwen3-8B}} \\
\midrule
 & 0.3 & KVCore & \textbf{29.73} & \textbf{44.20} & 53.95 & \textbf{62.62} & 48.75 & 33.79 & \textbf{33.70} & \textbf{24.45} & \textbf{24.23} & 90.26 & 40.24 & 39.00 & \textbf{9.50} & \textbf{94.54} & 66.78 & 62.25 & 47.37 \\
 &  & SnapKV & 28.92 & 44.17 & \textbf{55.27} & 62.05 & \textbf{49.58} & \textbf{35.88} & 33.10 & 24.19 & 24.21 & \textbf{90.46} & 39.89 & \textbf{41.50} & \textbf{9.50} & 92.43 & \textbf{67.12} & \textbf{62.52} & \textbf{47.55} \\
 &  & ChunkKV & 28.35 & 43.42 & 54.94 & 62.30 & 48.62 & 35.78 & 33.20 & 24.43 & 24.11 & 90.06 & \textbf{40.39} & 40.00 & 8.50 & 91.35 & 66.91 & 62.51 & 47.18 \\
\addlinespace[1pt]
 & 0.4 & KVCore & \textbf{30.79} & 44.50 & 52.90 & 62.28 & \textbf{50.61} & 35.00 & \textbf{33.03} & 23.86 & \textbf{23.93} & 90.16 & \textbf{40.90} & 35.50 & 9.00 & \textbf{95.79} & 66.85 & 62.88 & \textbf{47.37} \\
 &  & SnapKV & 29.03 & 43.88 & \textbf{54.56} & 62.05 & 48.88 & \textbf{35.57} & 32.78 & 24.34 & 23.78 & 89.96 & 40.68 & 40.00 & \textbf{9.50} & 91.77 & \textbf{67.39} & 62.72 & 47.31 \\
 &  & ChunkKV & 28.15 & \textbf{44.54} & 54.44 & \textbf{62.39} & 49.05 & 34.31 & 32.72 & \textbf{24.71} & 23.54 & \textbf{90.21} & 40.14 & \textbf{41.50} & 7.50 & 92.02 & 67.20 & \textbf{62.94} & 47.21 \\
\addlinespace[1pt]
 & 0.5 & KVCore & \textbf{29.42} & 43.58 & 52.20 & \textbf{62.36} & 49.60 & 33.77 & \textbf{32.85} & 23.73 & \textbf{23.48} & \textbf{90.24} & 40.24 & 34.00 & \textbf{9.50} & \textbf{97.08} & 65.70 & 62.85 & 46.91 \\
 &  & SnapKV & 29.07 & 43.16 & \textbf{54.87} & 61.63 & 49.28 & \textbf{35.90} & 32.25 & 24.19 & 23.21 & 89.96 & \textbf{40.62} & 39.50 & 8.50 & 92.39 & \textbf{68.15} & \textbf{63.36} & \textbf{47.25} \\
 &  & ChunkKV & 28.35 & \textbf{43.62} & 54.58 & 61.69 & \textbf{50.08} & 35.10 & 32.50 & \textbf{24.42} & 22.86 & 90.18 & 39.90 & \textbf{40.00} & 7.00 & 91.60 & 67.29 & 63.09 & 47.02 \\
\addlinespace[1pt]
 & 0.6 & KVCore & 27.99 & 42.84 & 52.30 & 62.11 & 47.76 & 33.09 & \textbf{32.54} & 23.39 & \textbf{22.92} & \textbf{90.24} & 39.82 & 32.00 & 8.05 & \textbf{96.85} & 66.22 & 63.89 & 46.38 \\
 &  & SnapKV & \textbf{28.61} & 43.17 & 53.94 & 61.64 & \textbf{49.78} & \textbf{35.83} & 31.79 & \textbf{24.28} & 22.43 & 90.21 & \textbf{40.85} & \textbf{38.00} & \textbf{9.00} & 93.04 & \textbf{67.49} & \textbf{64.21} & \textbf{47.14} \\
 &  & ChunkKV & 27.98 & \textbf{43.72} & \textbf{54.44} & \textbf{62.17} & 49.56 & 34.98 & 31.83 & 23.69 & 22.39 & 89.93 & 40.62 & 37.00 & 7.05 & 92.56 & 67.42 & 63.91 & 46.83 \\
\addlinespace[1pt]
 & 0.7 & KVCore & 29.47 & 42.39 & 52.54 & \textbf{62.18} & 47.12 & 32.46 & \textbf{31.68} & 23.30 & \textbf{21.78} & 88.92 & 39.88 & 29.50 & 7.00 & \textbf{97.10} & 65.77 & 64.20 & 45.96 \\
 &  & SnapKV & 28.12 & 42.19 & 52.04 & 62.13 & \textbf{49.41} & \textbf{35.58} & 30.73 & \textbf{24.26} & 21.35 & \textbf{90.21} & \textbf{41.14} & \textbf{34.75} & \textbf{7.55} & 93.10 & \textbf{67.80} & \textbf{64.93} & \textbf{46.58} \\
 &  & ChunkKV & \textbf{29.62} & \textbf{42.52} & \textbf{54.35} & 61.99 & 48.61 & 34.77 & 31.10 & 23.21 & 21.49 & \textbf{90.21} & 40.38 & 33.50 & 7.00 & 91.99 & 67.56 & 64.22 & 46.41 \\
\bottomrule
\end{tabular}%
}
\caption{Overall LongBench evaluation results of different KV cache compression methods under varying compression ratios. The table compares KVCore, SnapKV, and ChunkKV across three representative LLMs: Llama-3.1-8B-Instruct, Mistral-7B-Instruct-v0.3, and Qwen3-8B, with the goal of evaluating the impact of KV compression on downstream task accuracy. Columns are organized by LongBench task category, including single-document QA, multi-document QA, summarization, few-shot learning, synthetic tasks, and code completion. For each model block, we report per-task scores as well as the average score over all 16 LongBench tasks under different compression ratios. All methods compress only the prefill-stage KV cache, while the decode stage remains uncompressed. For each model and compression ratio, the best score among the three methods on each dataset and on the overall average metric is highlighted in bold. Higher values indicate better performance.}
\label{tab:longbench16-all-models}
\end{table*}


\textbf{Model generation quality.} \autoref{tab:longbench16-all-models} presents the accuracy of different KV cache compression methods across diverse LongBench task categories. The evaluation is conducted on three representative open-source LLMs under multiple compression ratios ranging from 0.3 to 0.7. The results demonstrate that KVCore consistently preserves generation quality across different models and datasets, exhibiting almost negligible accuracy degradation compared with SnapKV and ChunkKV, which represent classical token-wise and state-of-the-art block-wise KV compression approaches, respectively. Across all evaluation settings, the accuracy gap between KVCore and the best-performing configuration is typically within 1\%. Moreover, KVCore achieves the highest accuracy in several task and compression settings. These results indicate that substantial redundancy still exists in current KV cache representations, and that an appropriately designed block-level KV management and compression strategy can significantly reduce KV cache overhead while largely preserving model generation quality. As shown in \autoref{fig:ppl}, in long-context language modeling, KVCore does not noticeably amplify the perplexity variation caused by increasing context length. The PG19 results show that the PPL curve of KVCore remains within the same range as SnapKV and ChunkKV, without clear instability at longer context lengths. This suggests that KVCore preserves modeling quality comparable to representative baselines while retaining structured block-level information.

\begin{figure}[htbp]
    \centering
    \includegraphics[width=0.8\linewidth]{figure/pg19_dense_position_ppl_logy_max20k_paper_acm.pdf}
    \caption{Perplexity of different KV cache compression methods on PG19 as the context length increases. Experiments are conducted using Llama-3.1-8B-Instruct with a compression ratio of 0.5. The x-axis denotes the compressed context length, while the y-axis reports continuation perplexity (PPL) on the subsequent ground-truth text. For each evaluation point, PPL is computed over a 256-token continuation window.}
    \label{fig:ppl}
\end{figure}

\textbf{TTFT and throughout.} In order to compare the benefits of sparseness and offloading, we compared it with native vllm and infiniGen to compare the average latency of each token and the overall system throughput.At the same time, the request preemption rate is compared with the native vllm.Compared with the two baselines, kvcore reduces the average delay by 20\%, increases the system throughput by 1.5 times, and reduces the request preemption rate by 90\%, with almost no impact on accuracy.

\subsection{Sensitivity Study}
\textbf{compress ration.} KVCore

\textbf{layer reuse.} KVCore

\textbf{offload gb.} KVCore


\subsection{Overhead Analysis}

\begin{figure}
    \centering
    \includegraphics[width=\linewidth]{figure/sparse_index_overhead_paper_acm_wide.pdf}
    \caption{Sparse-index construction overhead comparison across SnapKV, ChunkKV, and KVCore. (a) Overhead under different request lengths. (b) Overhead under different batch sizes. (c) KVCore overhead under different reuse steps.}
    \label{fig:sparse-overhead}
\end{figure}

\textbf{Time Overhead.} \autoref{fig:sparse-overhead} compares the overhead of sparse index construction across different methods, including only importance score computation and top-k/index construction while excluding K/V gather overhead. SnapKV and ChunkKV both perform token-level scoring, causing their overhead to increase significantly with larger context lengths and batch sizes. ChunkKV further introduces additional chunk-level index construction, resulting in the highest overhead. In contrast, KVCore performs sparse selection at block granularity and reuses precomputed lightweight block summaries. As a result, online scoring and selection are performed over substantially fewer candidates, leading to lower and more scalable indexing overhead. With cross-step summary reuse, the amortized overhead of KVCore remains consistently low.

\textbf{Space Overhead.} The additional space overhead of KVCore mainly comes from reusable block summaries and block indices. By sharing and merging summaries across attention heads and layers, KVCore compresses summary metadata by nearly two orders of magnitude compared to the original KV cache, requiring only a small amount of persistent GPU memory.SnapKV and ChunkKV do not maintain additional summaries and therefore have lower static memory overhead. However, after KV offloading, they typically need to re-access large amounts of KV data to recompute token- or chunk-level importance scores, introducing additional data movement and online computation overhead.

\section{Related Work}
\label{sec:related_work}

\noindent\textbf{KV cache compression and sparse selection.}
Many KV cache compression methods exploit the non-uniform importance of historical KV entries to reduce attention computation and KV memory footprint.
Representative methods retain heavy-hitter tokens~\cite{zhang2023h2o}, compress prompt KV before generation~\cite{li2024snapkv}, allocate KV budgets across layers~\cite{cai2024pyramidkv,yang2024pyramidinfer}, perform query-aware sparse selection~\cite{tang2024quest}, group tokens into semantic chunks~\cite{liu2026chunkkv}, or prune KV cache in a structured manner~\cite{chitty2026pagedeviction,devoto2025expected,zhao2026hitkv,behnam2025rocketkv}.
These works mainly use sparsity as a runtime compression or per-step token selection rule to fit the KV cache into a restricted local memory budget.
In contrast, this work primarily treats online sparse scores as an indicative signal of the future memory access trajectory, shifting the focus from tensor-level pruning to runtime-level scheduling and lifecycle management.

\noindent\textbf{KV cache offloading and multi-tier management.}
Another line of work reduces GPU memory pressure by placing KV cache or inference states across heterogeneous memory tiers, including GPU memory, CPU memory, SSDs, and disaggregated storage.
These systems expand effective memory capacity, overlap data movement with computation, reuse prefix KV cache, or optimize prefill/decode execution~\cite{sheng2023flexgen,lee2024infinigen,xiong2024layerkv,xu2024pie,sun2025breaking,lin2026kvdrive,hu2026bidaw,qin2024mooncake,chen2025impress,liu2025lmcache,liu2024cachegen,yao2025cacheblend}.
They mainly address the multi-tier storage layout, underlying data movement primitives, and overlapping techniques.
The scheduling policies in these multi-tier management systems typically rely on static profiling, capacity-driven rules, or heuristic-based patterns such as LRU.
A dynamic, query-aware sparse scoring mechanism can serve as a complementary policy signal, providing a predictive guide for proactive offloading and prefetching.

\noindent\textbf{KV cache quantization.}
KV cache quantization reduces memory and bandwidth overhead by representing key and value states with fewer bits.
Prior work explores low-bit, asymmetric, and online vector quantization to support longer contexts and larger batches under limited memory capacity~\cite{liu2024kivi,hooper2024kvquant,zandieh2025turboquant}.
This line of work mainly optimizes the numerical representation of KV states rather than their access pattern.

\noindent\textbf{Efficient attention and LLM serving systems.}
Efficient attention methods reduce long-context inference cost by modifying attention patterns or model behavior, such as streaming attention with attention sinks and hardware-aligned sparse attention~\cite{xiao2024efficient,yuan2025native}.
LLM serving systems improve throughput and latency through continuous batching, block-based KV memory management, phase splitting, prefill/decode disaggregation, prefix reuse, and scheduling optimizations~\cite{kwon2023efficient,yu2022orca,agrawal2024taming,gim2024prompt,juravsky2024hydragen,jin2024p,patel2024splitwise}.
Under high GPU memory utilization, these systems traditionally rely on reactive and coarse-grained mechanisms, such as full-request preemption or full-sequence swapping, to handle budget exhaustion.



\section{Conclusion}
\label{sec:conclusion}
Existing KV sparsity and KV offloading systems are typically designed and optimized independently, overlooking the fact that sparsity itself implicitly reveals future KV access traces and can therefore guide runtime-level data movement and KV lifecycle management decisions. In this paper, we present KVCore, a KV-centric runtime for long-context LLM serving that jointly optimizes sparse computation and KV movement by elevating block-level sparsity into a runtime-level future-access signal.
KVCore introduces a lightweight block-level sparsity estimation mechanism to predict future KV importance with low overhead, and leverages this signal to jointly manage KV residency, offloading, prefetching, and eviction, thereby reducing redundant GPU–CPU transfers and unnecessary attention computation. Furthermore, KVCore exploits the similarity of sparse patterns across layers to reuse sparse indices and further amortize runtime sparsity planning overhead.
Experimental results show that KVCore effectively alleviates KV memory pressure and significantly reduces request-level preemption in long-context serving workloads. While maintaining near-lossless model accuracy, KVCore achieves up to 2.5× higher throughput compared with existing baselines.





%%
%% The next two lines define the bibliography style to be used, and
%% the bibliography file.
\bibliographystyle{ACM-Reference-Format}
\bibliography{samples/main}
\end{document}
\endinput
%%
%% End of file `sample-sigconf-authordraft.tex'.
